% ======================================================================
%  Hybrid-Sentinel IEEE Conference Paper (section-wise figures)
%  Every figure is introduced in body text: what it shows + what it means
% ======================================================================
\documentclass[10pt,conference]{IEEEtran}

\usepackage{amsmath,amssymb}
\usepackage{graphicx}
\usepackage{booktabs}
\usepackage{cite}
\usepackage{url}
\usepackage[hidelinks]{hyperref}
\usepackage{multirow}
\usepackage{caption}
\usepackage{subcaption}
\usepackage{placeins}   % \FloatBarrier keeps figures inside their section

\captionsetup{font=footnotesize, labelfont=bf, skip=4pt}
\captionsetup[sub]{font=scriptsize, labelfont=normalfont, skip=2pt}

\makeatletter
\setlength{\@fptop}{0pt}
\makeatother

\begin{document}

\title{Hybrid-Sentinel: A Three-Tier Confidence-Gated AI Firewall
       for NDN Edge Deployment}

\author{%
  \IEEEauthorblockN{Raju Guguloth}
  \IEEEauthorblockA{Indian Institute of Technology Madras\\
    M.Tech Research --- TCL Collaboration\\
    \texttt{[your.email@iitm.ac.in]}}
}

\maketitle

% ======================================================================
\begin{abstract}
Signature-based intrusion detection systems such as Snort remain the
backbone of network defense, yet their reliance on static,
hand-maintained rule sets leaves them blind to novel attacks, prone to
false alarms, and ill-suited to Named Data Networking (NDN) edge routers
that forward discrete Interest/Data units rather than IP flows.
This paper presents \textit{Hybrid-Sentinel}, a confidence-gated
behavioural firewall that inspects 20-packet windows of 17 features and
emits auditable \texttt{ALLOW}, \texttt{FLAG}, or \texttt{BLOCK}
decisions. We evaluate the same methodological pattern on two
complementary tracks. Track~A is a production three-tier cascade on
Docker IP lab traffic: a 156\,KB Random Forest gate, a
temperature-calibrated CNN-GRU (ONNX) over six attack classes, and a
Mahalanobis one-class detector for zero-day novelty---achieving 100\%
attack detection, 1.11\% benign false-positive rate, and 4.19\,ms mean
CPU latency on 14{,}219 held-out sequences. Track~B is an NDN forwarder
proof-of-concept that observes Pending Interest Table (PIT) and Content
Store (CS) state to detect Interest Flooding and Cache Pollution,
reaching 0.995 macro-F1 and 0.67\% benign FPR on 51{,}656 held-out
windows. Together, the results show that a staged, learning-based design
can outperform static signature engines on both IP-style behaviours and
NDN-native attacks, while remaining explainable for edge operation.
\end{abstract}

\begin{IEEEkeywords}
Named Data Networking, intrusion detection, Snort, cascade classifier,
Interest Flooding, Cache Pollution, ONNX, anomaly detection
\end{IEEEkeywords}

% ======================================================================
\section{Introduction}
\label{sec:intro}

Classical network intrusion detection is dominated by signature engines
such as Snort and Suricata. These systems are fast at matching known
patterns, but they fail when payloads are obfuscated, when attacks are
novel, or when the protocol surface changes---problems that become acute
at the NDN edge, where adversaries target forwarder state (Pending
Interest Table exhaustion, Content Store pollution) rather than TCP/IP
headers alone~\cite{ndn,afiha}. Deep learning detectors improve
coverage, yet running a heavyweight model on every packet rarely fits a
router latency budget.

Named Data Networking routes on content names and naturally inspects
discrete Interest and Data packets at each hop~\cite{ndn}. This
motivates a \emph{hybrid} firewall: a cheap gate on most traffic and a
deeper model only when suspicion rises~\cite{tcl-wiki}.
Hybrid-Sentinel implements that pattern and validates it on two tracks
without conflating their evidence:
\begin{itemize}
  \item \textbf{Track A (R18 production):} a three-tier cascade trained
        on controlled Docker TCP/IP PCAPs (brute force, HTTP flood, slow
        HTTP, port scan, DNS tunneling).
  \item \textbf{Track B (NDN PoC):} the same behavioural-window idea
        applied to NDN-native features extracted from a simulated PIT+CS
        forwarder under Interest Flooding and Cache Pollution.
\end{itemize}

\paragraph{Contributions.}
(1)~A confidence-gated three-tier cascade with shared feature contract,
ONNX deployment, and per-tier audit traces. (2)~Honest dual evaluation:
IP lab cascade metrics plus transferable NDN-native detection of Interest
Flooding and Cache Pollution. (3)~A reproducible build path from
PCAP/simulation to FastAPI service suitable for edge research deployment.

% ======================================================================
\section{Related Work}
\label{sec:related}

Snort-style IDS relies on deterministic pattern matching and continuous
rule maintenance; zero-day and protocol-native NDN attacks fall outside
typical IP signature libraries~\cite{afiha}. Flow-based CNN/RNN IDS on
public IP datasets (e.g., CIC-IDS) demonstrates deep models on packet
sequences~\cite{cicids} but does not address NDN forwarder state. Prior
NDN security work studies Interest Flooding against the PIT and
pollution against the CS~\cite{afiha,compagno}. Our contribution is a
\emph{confidence-gated cascade} for production IP behaviours, plus a
PoC showing that windowed behavioural features transfer to NDN state
observations---without claiming the IP models were trained on NDN
packets.

\FloatBarrier
% ======================================================================
\section{System Architecture}
\label{sec:system}

\subsection{Dual-Track Overview}

Fig.~\ref{fig:unified} is the master map of Hybrid-Sentinel. Reading
top to bottom, it conveys four messages: (1)~NDN edge routers inspect
Interest/Data units and maintain PIT/CS state, which signature engines
such as Snort were not designed to protect; (2)~the design choice is a
\emph{cheap gate on most traffic} and deeper ML only on suspicion, with
auditable \texttt{ALLOW}/\texttt{FLAG}/\texttt{BLOCK} actions; (3)~both
tracks share the same methodological contract---17 behavioural features
organised as 20-packet windows; (4)~Track~A (left) is the production
IP cascade (Random Forest $\rightarrow$ CNN-GRU $\rightarrow$
Mahalanobis), while Track~B (right) instantiates the same window idea
on NDN-native PIT/CS features for Interest Flooding and Cache Pollution.
The figure therefore makes the honesty claim explicit: the
\emph{method} is shared; the \emph{trained models and feature semantics}
are track-specific.

\begin{figure*}[!t]
  \centering
  \includegraphics[width=\textwidth]{figures/fig_unified_system_architecture.pdf}
  \caption{Dual-track Hybrid-Sentinel architecture (full page width).
    Top: NDN motivation vs.\ Snort limitations and hybrid design choice.
    Middle: shared 17-feature / 20-packet detection contract.
    Bottom-left: Track~A R18 cascade on IP lab traffic.
    Bottom-right: Track~B NDN forwarder PoC on PIT/CS state.}
  \label{fig:unified}
\end{figure*}

\subsection{Track A: Runtime Cascade and Build Pipeline}

Each Track~A decision uses a scaled window
$\mathbf{X}\in\mathbb{R}^{20\times 17}$. Tier-1 estimates
$P(\text{BENIGN})$; above $\tau_1{=}0.90$ the window is \texttt{ALLOW}ed
without deep inference. Otherwise Tier-2 (CNN-GRU, temperature-scaled)
assigns one of six labels; confidence above $\tau_b{=}0.95$ yields
\texttt{BLOCK}, and $[\tau_f,\tau_b)$ with $\tau_f{=}0.80$ yields
\texttt{FLAG}. Tier-3 scores Mahalanobis distance on a 128-D embedding
and may \texttt{FLAG} residual novelty as \texttt{ANOMALY}.

Fig.~\ref{fig:arch} conveys the \emph{runtime} decision path for one
window. The top box is the input contract. The diamond after Tier-1 is
the confidence gate: high benign probability exits immediately as
\texttt{ALLOW} (green), which keeps clean traffic cheap. The orange
escalate path enters Tier-2, where calibrated class probabilities become
operational actions. Tier-3 only acts on residual uncertainty, flagging
embedding outliers without inventing an attack class name. The bottom
box emphasises that every path ends in one of three auditable actions
for SOC review.

\begin{figure}[!t]
  \centering
  \includegraphics[width=\columnwidth]{figures/architecture_flow.pdf}
  \caption{Track~A runtime cascade. What it conveys: Tier-1 fast
    \texttt{ALLOW} when $P(\text{BENIGN})\ge 0.90$; escalated windows
    reach Tier-2 (six-class ONNX) then optional Tier-3 novelty
    \texttt{FLAG}; final action is always \texttt{ALLOW}, \texttt{FLAG},
    or \texttt{BLOCK}.}
  \label{fig:arch}
\end{figure}

Fig.~\ref{fig:orch} conveys the \emph{build} story, not packet scoring.
Left-to-right stages show PCAP capture $\rightarrow$ 17-feature
extraction $\rightarrow$ 20-packet windows $\rightarrow$ three-tier
training. The middle row shows ONNX export, Tier-3 calibration, and the
single cascade module shared by offline evaluation and the live API.
The bottom row mirrors runtime tiers and final actions. This figure is
the reproducibility map: any retrain must regenerate the same export
chain so evaluation and serving cannot drift.

\begin{figure}[!t]
  \centering
  \includegraphics[width=\columnwidth]{figures/orchestration_pipeline.pdf}
  \caption{Track~A build orchestration. What it conveys: how lab PCAPs
    become exported ONNX/cascade artifacts and a FastAPI service with
    one shared decision logic.}
  \label{fig:orch}
\end{figure}

\subsection{Track B: NDN Forwarder Topology}

Fig.~\ref{fig:ndn_topo} conveys the NDN experiment setup. Green benign
consumers issue Zipf-popular, satisfiable Interests. Red attackers inject
either unsatisfiable unique names (Interest Flooding $\rightarrow$ PIT
exhaustion) or cold but valid names (Cache Pollution $\rightarrow$ CS
hit-ratio collapse). Router~$R$ is the observation point: PIT/CS
counters, name entropy, and unsatisfied-Interest ratios are logged
causally. Interests leave~$R$ toward the producer and Data returns on
the reverse path. The footer states the PoC pipeline---17 NDN-native
features $\rightarrow$ 20-packet windows $\rightarrow$ Random Forest---
and reminds the reader this is a discrete-event simulation, not a live
NFD capture.

\begin{figure}[!t]
  \centering
  \includegraphics[width=\columnwidth]{figures/fig_ndn_architecture.pdf}
  \caption{Track~B NDN topology. What it conveys: where Interests/Data
    flow, why router~$R$ (PIT+CS) is monitored, and how windows feed the
    NDN classifier (simulation, not live NFD).}
  \label{fig:ndn_topo}
\end{figure}

\FloatBarrier
% ======================================================================
\section{Methodology}
\label{sec:method}

\subsection{Track A Features and Training}

Seventeen packet-level IP features (length, protocol flags, payload
entropy, flow aggregates) are extracted with Scapy. Sliding windows of
size 20 (stride 10) form sequences. Sequences are grouped by
pseudo-flow blocks; train/validation/test splits are 70/15/15 with zero
group overlap. The scaler is fit on training data only. Tier-2 uses
focal loss with elevated benign weight and temperature calibration on
validation. Tier-3 is fit on benign training embeddings and must be
re-exported when Tier-2 weights change. Tier-2 inference uses ONNX
Runtime on CPU; cascade logic lives in a single module shared by
evaluation and the API.

\subsection{Track B NDN Features}

NDN features deliberately parallel the production 17-dimensional
contract but read forwarder state: Interest/Data indicators, name depth
and entropy, PIT occupancy and growth, CS size and hit signals,
interarrival and windowed rates. Features are causal (past observations
only), so the extractor is conceptually deployable on a live forwarder.
A Random Forest classifier is trained with a stratified 75/25
train/test split on simulated windows.

\FloatBarrier
% ======================================================================
\section{Experimental Setup}
\label{sec:setup}

\paragraph{Track A lab topology.}
Traffic was generated in an isolated Docker bridge network with producer
containers (HTTP, SSH, DNS) and scripted attackers. Fig.~\ref{fig:docker}
conveys the physical capture layout used for Track~A: clients and
attackers share a bridge with service producers so every PCAP is labelled
by ground-truth scenario (benign baseline vs.\ one attack type). The
held-out test set contains 14{,}219 sequences (814 benign, 13{,}405
attack) never used for training or model selection. Latency is measured
over 2{,}000 CPU inferences (single-threaded Tier-1).

\begin{figure}[!t]
  \centering
  \includegraphics[width=\columnwidth]{figures/docker_topology.pdf}
  \caption{Track~A Docker lab topology. What it conveys: how labelled
    IP PCAPs are produced (benign services + scripted attacks) before
    feature extraction and cascade evaluation.}
  \label{fig:docker}
\end{figure}

Fig.~\ref{fig:classdist} conveys class imbalance on the Track~A held-out
set. Attack classes dominate sample counts; benign is the minority class
that false-positive rate cares about. This imbalance is why Tier-2 uses
focal loss and elevated benign weight, and why E2E reporting separates
attack detection from benign FPR.

\begin{figure}[!t]
  \centering
  \includegraphics[width=\columnwidth]{figures/class_distribution.pdf}
  \caption{Track~A held-out class distribution. What it conveys: attack
    classes outnumber benign, so metrics must not hide benign false
    alarms behind overall accuracy.}
  \label{fig:classdist}
\end{figure}

\paragraph{Track B simulation.}
A discrete-event NDN simulator emits 206{,}624 windows (43{,}299
BENIGN, 63{,}308 INTEREST\_FLOODING, 100{,}017 CACHE\_POLLUTION).
Held-out evaluation uses 51{,}656 windows. Absolute numbers reflect
simulator assumptions (Poisson arrivals, Zipf popularity, LRU CS, fixed
PIT lifetime) and complement---not replace---Track~A.

\FloatBarrier
% ======================================================================
\section{Results}
\label{sec:results}

\subsection{Track A: Per-Tier Discrimination}

Table~\ref{tab:tier} lists per-component metrics. Tier-1 separates
benign fast-path from attack escalation with zero benign false
escalations on the test set. Tier-2 macro-F1 is 0.985;
\texttt{BRUTE\_FORCE} has the lowest recall (0.914) due to early-packet
overlap with slow HTTP / flood patterns.

\begin{table}[!t]
  \caption{Track~A (R18) per-tier metrics on held-out test
    ($n = 14{,}219$). Tier-2 uses macro averaging.}
  \label{tab:tier}
  \centering\small
  \begin{tabular}{lccccc}
    \toprule
    \textbf{Component} & \textbf{Acc.} & \textbf{Prec.} &
    \textbf{Rec.} & \textbf{F1} & \textbf{Latency} \\
    \midrule
    Tier-1 gate    & 1.000 & 1.000 & 1.000 & 1.000 & 3.66\,ms \\
    Tier-2 CNN-GRU & 0.981 & 0.986 & 0.985 & 0.985 & 0.18\,ms \\
    Tier-3 anomaly & 0.995 & 0.992 & 1.000 & 0.996 & 0.21\,ms \\
    Full cascade   & 0.999 & 0.999 & 1.000 & 1.000 & 4.19\,ms \\
    \bottomrule
  \end{tabular}

  \parbox{\columnwidth}{\footnotesize
    Throughput: Tier-1 272\,seq/s; Tier-2 ONNX 5{,}570\,seq/s; full
    cascade 238\,seq/s. Baseline gradient-boosting latency
    $\approx$13\,ms.}
\end{table}

Fig.~\ref{fig:metrics} conveys whether each tier is individually
competent. Bars for accuracy, precision, recall, and F1 stay near 1.0
across Tier-1, Tier-2, Tier-3, and the full cascade, showing that the
cascade does not trade security for speed---it composes strong
components.

\begin{figure}[!t]
  \centering
  \includegraphics[width=\columnwidth]{figures/tier_metrics_bars.pdf}
  \caption{Track~A tier metric comparison. What it conveys: each stage
    and the full cascade remain high on Acc/Prec/Rec/F1; the cascade is
    not a weak link.}
  \label{fig:metrics}
\end{figure}

Fig.~\ref{fig:t1cm} conveys the Tier-1 gate task in binary form: true
benign windows should be allowed (fast path) and true attacks should be
escalated. Perfect separation on the held-out set means Tier-1 does not
leak attacks into silent \texttt{ALLOW}.

\begin{figure}[!t]
  \centering
  \includegraphics[width=\columnwidth]{figures/tier1_confusion.pdf}
  \caption{Tier-1 gate confusion (allow vs.\ escalate). What it conveys:
    on the held-out set, benign windows are not falsely escalated and
    attacks are not falsely allowed.}
  \label{fig:t1cm}
\end{figure}

Fig.~\ref{fig:cm} conveys Tier-2 class confusion. The dark diagonal
means near-perfect recall for five classes. The only material
off-diagonal mass is
\texttt{BRUTE\_FORCE}$\rightarrow$\texttt{DDOS\_HTTP} ($\approx$0.09):
early packets of login-brute sequences can look flood-like, so a small
fraction is typed as HTTP flood rather than brute force. This is a
\emph{typing} error among attacks, not a miss of attack presence.

\begin{figure}[!t]
  \centering
  \includegraphics[width=\columnwidth]{figures/confusion_matrix.pdf}
  \caption{Tier-2 normalised confusion matrix. What it conveys: five
    classes are nearly perfect; the residual error is
    BRUTE\_FORCE$\leftrightarrow$DDOS\_HTTP overlap ($\approx$9\%).}
  \label{fig:cm}
\end{figure}

Fig.~\ref{fig:heatmap} conveys the same Tier-2 story as per-class
precision, recall, and F1 so the weak class is visible without reading
matrix cells. \texttt{BRUTE\_FORCE} recall is the lowest bar; other
classes sit at or near 1.0.

\begin{figure}[!t]
  \centering
  \includegraphics[width=\columnwidth]{figures/per_class_heatmap.pdf}
  \caption{Tier-2 per-class precision/recall/F1. What it conveys:
    BRUTE\_FORCE is the weakest class; DNS tunnelling, port scan, and
    slow HTTP are essentially saturated.}
  \label{fig:heatmap}
\end{figure}

\subsection{Track A: Cascade Load, Security Outcome, and Latency}

Fig.~\ref{fig:funnel} conveys \emph{where traffic goes} in the cascade.
Only a minority of windows exit at Tier-1; most escalate to Tier-2; a
smaller fraction reach Tier-3. That shape is intentional: the expensive
models run only when suspicion rises.

\begin{figure}[!t]
  \centering
  \includegraphics[width=\columnwidth]{figures/cascade_funnel.pdf}
  \caption{Cascade funnel (\% of sequences per stage). What it conveys:
    cheap-first behaviour---most suspicious mass reaches Tier-2; Tier-3
    sees only residual uncertainty.}
  \label{fig:funnel}
\end{figure}

Fig.~\ref{fig:e2e} conveys the operational security headline for
Track~A: attack detection rate, benign false-positive rate, and
end-to-end precision/F1. The full cascade flags or blocks all 13{,}405
attacks (100\%) and falsely flags 9 of 814 benign windows (1.11\% FPR);
none are auto-blocked without review.

\begin{figure}[!t]
  \centering
  \includegraphics[width=\columnwidth]{figures/e2e_security.pdf}
  \caption{Track~A end-to-end security rates. What it conveys: 100\%
    attack detection with 1.11\% benign FPR on the held-out set---the
    production cascade metric that matters for SOC deployment.}
  \label{fig:e2e}
\end{figure}

Fig.~\ref{fig:t3roc} conveys Tier-3 ranking quality for novelty
detection. An ROC curve near the top-left with AUC\,$=$\,1.0 on held-out
scores means Mahalanobis distance separates benign embeddings from
attack/novel windows under the lab distribution (with the usual caveat
that lab AUC can be optimistic for open-world traffic).

\begin{figure}[!t]
  \centering
  \includegraphics[width=\columnwidth]{figures/tier3_roc.pdf}
  \caption{Tier-3 Mahalanobis ROC. What it conveys: embedding novelty
    scores rank attack/benign windows almost perfectly on the lab test
    set (AUC\,$=$\,1.0).}
  \label{fig:t3roc}
\end{figure}

Fig.~\ref{fig:lat} conveys edge feasibility. Mean and p99 latency for
Tier-1, Tier-2 ONNX, Tier-3, and the full cascade are plotted against a
13\,ms gradient-boosting reference; the lower panel shows throughput
(seq/s). The full cascade averages 4.19\,ms---well below the baseline---
while Tier-1 dominates runtime cost, confirming the gate is the
throughput bottleneck and ONNX Tier-2 is cheap once escalated.

\begin{figure}[!t]
  \centering
  \includegraphics[width=\columnwidth]{figures/latency_throughput.pdf}
  \caption{Track~A latency and throughput (CPU, $n{=}2000$). What it
    conveys: full cascade 4.19\,ms mean (p99 6.49\,ms) vs.\ 13\,ms
    baseline; Tier-1 is the latency/throughput bottleneck.}
  \label{fig:lat}
\end{figure}

\FloatBarrier
\subsection{Track B: NDN PoC Discrimination}

Table~\ref{tab:ndn} summarises held-out NDN metrics. Interest Flooding
is perfectly separable (PIT-exhaustion signature). Residual errors are a
small BENIGN~$\leftrightarrow$~CACHE\_POLLUTION overlap during
low-density attack warmup---expected under Zipf/LRU dynamics.

\begin{table}[!t]
  \caption{Track~B NDN PoC held-out metrics ($n = 51{,}656$).}
  \label{tab:ndn}
  \centering\small
  \begin{tabular}{lcccc}
    \toprule
    \textbf{Class / Aggregate} & \textbf{Prec.} & \textbf{Rec.} &
    \textbf{F1} & \textbf{Acc.} \\
    \midrule
    BENIGN             & 0.987 & 0.993 & 0.990 & 0.993 \\
    CACHE\_POLLUTION   & 0.997 & 0.994 & 0.996 & 0.994 \\
    INTEREST\_FLOODING & 1.000 & 1.000 & 1.000 & 1.000 \\
    \midrule
    Overall            & ---   & ---   & 0.995 & 0.996 \\
    \bottomrule
  \end{tabular}

  \parbox{\columnwidth}{\footnotesize
    Attack detection 99.65\%; benign FPR 0.67\%;
    ROC-AUC (attack vs.\ benign) 0.9997.}
\end{table}

Fig.~\ref{fig:ndn_cm} conveys class-wise NDN errors. Interest Flooding
is a perfect diagonal (1.00 / 15{,}827). The only non-zero off-diagonal
cells are BENIGN$\rightarrow$CACHE\_POLLUTION (73 windows, 0.01) and
CACHE\_POLLUTION$\rightarrow$BENIGN (142 windows, 0.01). That pattern
means the PoC almost never confuses flooding with pollution; the hard
case is mild pollution warmup looking locally benign (and vice versa).

\begin{figure}[!t]
  \centering
  \includegraphics[width=\columnwidth]{figures/fig_ndn_confusion_matrix.pdf}
  \caption{NDN normalised confusion matrix. What it conveys: Interest
    Flooding is perfectly classified; residual errors are a 1\% swap
    between BENIGN and CACHE\_POLLUTION only.}
  \label{fig:ndn_cm}
\end{figure}

Fig.~\ref{fig:ndn_metrics} conveys the same result as precision, recall,
and F1 bars. All three classes sit near 1.0; Interest Flooding saturates
every bar, matching the PIT-exhaustion signature being unambiguous.

\begin{figure}[!t]
  \centering
  \includegraphics[width=\columnwidth]{figures/fig_ndn_per_class_metrics.pdf}
  \caption{NDN per-class precision/recall/F1. What it conveys: no weak
    NDN class; Flooding is saturated; Pollution and Benign remain
    $\ge$0.99 F1.}
  \label{fig:ndn_metrics}
\end{figure}

Fig.~\ref{fig:ndn_acc} conveys per-class accuracy (row-correct rate). It
is included so readers who prefer accuracy bars over P/R/F1 can confirm
the same ranking: Flooding 1.0, Pollution $\approx$0.994, Benign
$\approx$0.993.

\begin{figure}[!t]
  \centering
  \includegraphics[width=\columnwidth]{figures/fig_ndn_per_class_accuracy.pdf}
  \caption{NDN per-class accuracy. What it conveys: accuracy mirrors the
    confusion-matrix diagonal---Flooding perfect, Pollution/Benign
    slightly below due to warmup overlap.}
  \label{fig:ndn_acc}
\end{figure}

Fig.~\ref{fig:ndn_summary} conveys the four NDN headline rates in one
chart: accuracy, macro-F1, attack detection, and benign FPR. Detection
stays high while benign FPR stays at 0.67\%---the PoC analogue of
Track~A's E2E security panel.

\begin{figure}[!t]
  \centering
  \includegraphics[width=\columnwidth]{figures/fig_ndn_summary_metrics.pdf}
  \caption{NDN summary metrics. What it conveys: 99.58\% accuracy,
    0.995 macro-F1, 99.65\% attack detection, and 0.67\% benign FPR on
    51{,}656 held-out windows.}
  \label{fig:ndn_summary}
\end{figure}

Fig.~\ref{fig:ndn_roc} conveys ranking quality, not just hard decisions.
The left panel is binary attack-vs-benign ROC (AUC\,$=$\,0.9997); the
right panel is one-vs-rest ROC for each class (macro AUC\,$=$\,0.9998).
Curves hugging the top-left corner mean score thresholds remain stable
even when class priors differ.

\begin{figure}[!t]
  \centering
  \includegraphics[width=\columnwidth]{figures/fig_ndn_roc_curves.pdf}
  \caption{NDN ROC curves. What it conveys: near-perfect ranking for
    attack detection (left) and per-class one-vs-rest scores (right);
    thresholds are not brittle.}
  \label{fig:ndn_roc}
\end{figure}

\FloatBarrier
% ======================================================================
\section{Discussion}
\label{sec:discuss}

Relative to Snort-like static rules, Hybrid-Sentinel learns behavioural
windows and can FLAG novelty without waiting for a new signature.
Track~A shows that cascading a cheap gate with calibrated deep
classification and embedding anomaly detection meets edge latency
budgets on CPU (Figs.~\ref{fig:funnel}--\ref{fig:lat}). Track~B shows
the \emph{methodology} transfers when features are redefined for PIT/CS
state (Figs.~\ref{fig:ndn_topo}--\ref{fig:ndn_roc}).

\paragraph{Limitations.}
Track~B is a simulation, not a live NFD/ndnSIM capture. Track~A remains
IP-overlay lab traffic; enterprise TLS and live NDN malware-in-Data
remain future work. The paper deliberately separates the two tracks to
avoid overclaiming that production cascade weights were trained on NDN
packets.

% ======================================================================
\section{Conclusion}
\label{sec:concl}

Hybrid-Sentinel delivers a reproducible, confidence-gated AI firewall:
100\% attack detection with 1.11\% benign FPR and 4.19\,ms cascade
latency on IP lab traffic, plus an NDN PoC that detects Interest
Flooding and Cache Pollution at 0.995 macro-F1. By pairing a
Snort-motivated critique with dual-track evidence, we show that hybrid
behavioural detection is a practical path toward NDN edge defense.

\section*{Acknowledgment}
The author thanks the TCL--IITM research collaboration for
infrastructure and technical guidance.

\begin{thebibliography}{00}
\bibitem{tcl-wiki}
T.~Sood \textit{et al.},
``Next-Generation Firewall Architectures: Rule-Based vs.\ AI
Transformers,'' TCL--IITM GitLab Wiki, 2026.

\bibitem{ndn}
L.~Zhang \textit{et al.},
``Named Data Networking,''
\emph{ACM SIGCOMM CCR}, vol.~44, no.~3, pp.~66--73, 2014.

\bibitem{afiha}
A.~Compagno \textit{et al.},
``Poseidon: Mitigating Interest Flooding DDoS Attacks in Named Data
Networking,'' \emph{Proc.\ IEEE LCN}, 2013.

\bibitem{compagno}
A.~Compagno \textit{et al.},
``Botcache: Misusing Content Caching for Hidden Communication in
Named-Data Networking,'' \emph{Proc.\ ACM SIGCOMM ANCS}, 2013.

\bibitem{cicids}
I.~Sharafaldin \textit{et al.},
``Toward Generating a New Intrusion Detection Dataset and Intrusion
Traffic Characterization,'' \emph{Proc.\ ICISSp}, 2018.
\end{thebibliography}

\end{document}
